% Part: sets-functions-relations
% Chapter: sets
% Section: important-sets

\documentclass[../../../include/open-logic-section]{subfiles}

\begin{document}

\olfileid{sfr}{set}{imp}
\olsection{কয়েকটি গুরুত্বপূর্ণ সেট}

\begin{ex}
আমরা মূলত এমন সেট নিয়ে আলোচনা করব, যাদের !!{element}s গাণিতিক বস্তু।
এমন চারটি সেট এতটাই গুরুত্বপূর্ণ যে, তাদের নির্দিষ্ট নাম আছে:
\begin{multline*}
    \Nat = \{0, 1, 2, 3, \ldots\} \\
    \shoveright{\text{স্বাভাবিক সংখ্যার সেট}}\\
    \shoveleft{\Int = \{\ldots, -2, -1, 0, 1, 2, \ldots\}} \\
    \shoveright{\text{পূর্ণসংখ্যার সেট}}\\
    \shoveleft{\Rat = \Setabs{\nicefrac{m}{n}}{m, n \in \Int\text{ এবং }n \neq 0}}\\
    \shoveright{\text{মূলদ সংখ্যার সেট}}\\
    \shoveleft{\Real = (-\infty, \infty)}\\
    \text{বাস্তব সংখ্যার সেট (সতত সমষ্টি)}
\end{multline*}
এগুলি সবই \emph{অসীম} সেট; অর্থাৎ, প্রতিটিরই অসীমসংখ্যক
!!{element}s আছে।

এই সেটগুলির মধ্য দিয়ে এগোতে এগোতে আমরা সংখ্যার ভাণ্ডারে
\emph{আরও} সংখ্যা যোগ করছি। স্পষ্টতই,
$\Nat \subseteq \Int \subseteq \Rat \subseteq \Real$: প্রতিটি স্বাভাবিক
সংখ্যাই পূর্ণসংখ্যা; প্রতিটি পূর্ণসংখ্যাই মূলদ; আর প্রতিটি মূলদ সংখ্যাই বাস্তব।
একইভাবে, $\Nat \subsetneq \Int \subsetneq \Rat$ হওয়াও স্পষ্ট:
$-1$ পূর্ণসংখ্যা হলেও স্বাভাবিক সংখ্যা নয়, আর $\nicefrac{1}{2}$
মূলদ হলেও পূর্ণসংখ্যা নয়। তবে $\Rat \subsetneq \Real$, অর্থাৎ কিছু
বাস্তব সংখ্যা যে মূলদ নয়, তা অতটা স্পষ্ট নয়\oliflabeldef{sfr:arith:real:realline}{;
\olref[arith][real]{realline}-এ আমরা এই প্রসঙ্গে ফিরব}{}।

কখনও কখনও ধনাত্মক পূর্ণসংখ্যার সেট $\PosInt = \{1, 2, 3, \dots\}$
এবং শুধু প্রথম দুটি স্বাভাবিক সংখ্যা নিয়ে গঠিত সেট $\Bin = \{0, 1\}$-ও
ব্যবহার করব।
\end{ex}

\begin{tagblock}{compsci}
\begin{ex}[প্রতীকক্রম]
আরও একটি আকর্ষণীয় উদাহরণ হল কোনও বর্ণমালা $A$-এর উপর গঠিত
\emph{সসীম প্রতীকক্রমগুলির} সেট $A^{*}$: $A$-এর উপাদানগুলির যে-কোনও
সসীম অনুক্রমই $A$-এর উপর একটি প্রতীকক্রম। প্রতিটি বর্ণমালা~$A$-এর জন্য,
$A$-এর উপর প্রতীকক্রমগুলির মধ্যে \emph{শূন্য প্রতীকক্রম $\Lambda$}-কেও
অন্তর্ভুক্ত করি। যেমন,
\begin{multline*}
\Bin^*
=\{\Lambda,0,1,00,01,10,11,\\
000,001,010,011,100,101,110,111,0000,\ldots\}.
\end{multline*}
$x=x_{1}\ldots x_{n}\in A^{*}$ যদি $A$-এর $n$টি ``বর্ণ'' নিয়ে গঠিত
প্রতীকক্রম হয়, তবে তার \emph{দৈর্ঘ্য}~$n$ বলি এবং $\len{x}=n$ লিখি।
\end{ex}
\end{tagblock}

\begin{ex}[অসীম অনুক্রম]
যে-কোনও সেট $A$-এর জন্য, $A$-এর !!{element}sের অসীম অনুক্রমগুলির
সেট~$A^\omega$-ও বিবেচনা করতে পারি। একটি অসীম অনুক্রম
$a_1a_2a_3a_4\dots$ হল বস্তুর এমন একটি তালিকা, যা এক দিকে অসীম পর্যন্ত
চলে এবং যার প্রতিটি বস্তু $A$-এর !!a{element}।
\end{ex}

\end{document}
